% ============================================================================
%  Capítulo 7 — Discusión
% ============================================================================
\chapter{Discusión}
\label{cap:discusion}

\section{Interpretación de los hallazgos}
\label{sec:hipotesis_discusion}

\paragraph{Dinámica de la planta.} Al variar el $b_0$ del controlador entre $1/6$ y $6$ veces el
nominal, se evidencia que el lazo de control es capaz de soportar las discrepancias entre las
ganancias (Figuras~\ref{fig:adrc_variacion}(d) y \ref{fig:sat_adrc_gpi_variacion}(d)) sin entrar en
inestabilidad. Por lo que, incluso sin una ganancia exacta, el sistema de control por sí solo tiene
la suficiente robustez para seguir funcionando a pesar de las discrepancias que haya con la planta.
Ahora bien, esto es sin retardo, donde el margen de ganancia es infinito
(Figura~\ref{fig:adrc_rlocus_b0}(a)). Con $\theta = 50$~s el margen de ganancia es de $2{,}12$
(Figura~\ref{fig:adrc_margenes_retardo}), así que el lazo se vuelve inestable si el $b_0$ del
controlador queda por debajo de $b_{0,\mathrm{planta}}/2{,}12$ (Figura~\ref{fig:adrc_rlocus_b0}(b)). Por
eso, si hay duda, conviene que $b_0$ quede por encima del de la planta: el lazo pierde velocidad, pero no
robustez.

\paragraph{Retardo.} El sistema es capaz de aguantar las variaciones del retardo de la planta,
incluso hasta el doble del de diseño, $\theta = 100$~s, sin desestabilizarse
(Figura~\ref{fig:sat_adrc_gpi_variacion}(c)). La cuestión es que se ve un incremento en el sobrepico
y mayores oscilaciones ante una referencia tipo paso: con $\theta = 50$~s el margen de fase baja de
$65{,}8^{\circ}$ a $36{,}0^{\circ}$ (Figura~\ref{fig:adrc_margenes_retardo}), y recordemos que el
margen de fase se puede relacionar con el comportamiento transitorio de un sistema, más
específicamente con su coeficiente de amortiguamiento. Pero en las
Figuras~\ref{fig:adrc_bezier_retardo} y \ref{fig:pref_escalon}, con Bézier y con prefiltro, se observa
cómo el sobrepico casi desaparece: con el prefiltro y $\theta = 50$~s baja de $46$ a $0~\%$, y con
Bézier las respuestas de $\theta = 20$ a $60$~s prácticamente se superponen. El margen de fase del
lazo no cambia, lo que cambia es que la referencia ya no excita esa oscilación, por lo que el sistema
puede responder mucho mejor ante estas discrepancias con los datos de la planta, específicamente con
el retardo.

Ahora bien, para el retardo no sirve aumentar la relación $\omega_o/\omega_c$ subiendo $\omega_o$. En la
Figura~\ref{fig:adrc_margen_fase_kESO}(a) y (c) el margen de fase sube un poco con esa relación, pero eso es
sin retardo. El cruce de ganancia se mueve más rápido, de $37{,}5$ a $66{,}9$~rad/h con
$\omega_o = 10\,\omega_c$, por lo que el retardo máximo \eqref{eq:mapa_robustez} baja de $110$ a
$65$~s, y con $\omega_o/\omega_c \approx 14$ ya queda por debajo de los $50$~s
(Tabla~\ref{tab:adrc_frecuencia_kESO}). Es decir, si el retardo empieza a afectar el sistema, lo que
hay que hacer es bajar $\omega_o$, no subirlo: con $\omega_o/2$ ($\omega_o/\omega_c = 2{,}5$) se llega a $180$~s. Bajar
$\omega_c$ también sirve, porque el cruce casi no se mueve: con $\omega_c/4$ se llega a $154$~s
(Tabla~\ref{tab:adrc_frecuencia_kESO}).
La otra opción, si el retardo se conoce, es sincronizar el observador \eqref{eq:disc_retardo}: con
$\theta = 50$~s el pico de oxígeno baja y se comienza a comportar mejor. 

\paragraph{Manejo de la saturación.} Los esquemas de manejo de la saturación, el del ESO
(Figura~\ref{fig:simulink_eso}, con resultados en la Figura~\ref{fig:adrc_lazo_cerrado_saturacion}) y
el del GPI (Figura~\ref{fig:simulink_gpi}, con resultados en la Figura~\ref{fig:sat_arranque}),
muestran muy buenos resultados: prácticamente logran manejar la saturación y el lazo no se
desestabiliza, incluso ante variaciones de los parámetros de la planta
(Figura~\ref{fig:sat_adrc_gpi_variacion}). Cuando el sistema se satura y regresa no se observa
sobrepico y se mantiene en el régimen de operación presentado en el diseño. En caso contrario, cuando
no hay manejo de saturación, el controlador se demora más en responder después de salir de saturación
y muestra sobrepicos más grandes: en la Figura~\ref{fig:sat_arranque} el GPI sin \emph{anti-windup} se
queda unos $16$~min en $100~\%$ y el oxígeno sube hasta $2{,}79$~mg/L. Lo que ningún esquema evita es
que, mientras el soplador está en el límite, el oxígeno suba más despacio que en el diseño: tarda
$354$~s en alcanzar el $63~\%$ del escalón, frente a $98$~s sin saturación.

\paragraph{Manejo del ruido.} Para estos controladores con enfoque de rechazo de perturbaciones se
tiene un manejo del ruido muy adecuado, ya que, como se ha venido mencionando, la función de
transferencia del ADRC basado en ESO es estrictamente propia \eqref{eq:gpi_pi_filtro}. Como se mostró
en la Sección~\ref{sec:impl_ruido}, con \eqref{eq:ruido_ft}, se pueden ajustar las frecuencias que se
desean filtrar y las que se desean estimar a partir de las ganancias del controlador, más
específicamente del ancho de banda del observador. Por lo que el ruido, junto con el retardo, es
quien realmente pone la cota superior a las ganancias del observador. Además, en la
Figura~\ref{fig:adrc_gang_kESO}(h) se ve la función de transferencia del ruido a $u$, y podemos
manejar qué tanto deja pasar el controlador al variar la relación $\omega_o/\omega_c$: entre más
grande, más ruido dejamos pasar, $10{,}5$~dB en alta frecuencia con $\omega_o = 5\,\omega_c$ y
$21{,}2$~dB con $10\,\omega_c$ (Tabla~\ref{tab:adrc_frecuencia_kESO}). Por eso el ancho de banda del
observador tampoco lo debemos dejar tan grande.

\paragraph{Discretización.} Con $T_s = 50$~ms todos los métodos, ya sea ZOH, Tustin, Euler hacia
adelante o Euler hacia atrás, dan una muy buena discretización
(Figura~\ref{fig:disc_metodos_bode}), porque la frecuencia de muestreo está muy por encima de la del
lazo cerrado. A medida que esa frecuencia baja, los métodos sí se empiezan a separar: con $10$ veces
el ancho de banda del lazo cerrado, $T_s = 40$~s, las magnitudes siguen juntas, pero en la fase, en el
cruce de ganancia, Tustin se separa del continuo $0{,}13^{\circ}$, Euler $1^{\circ}$ a $3^{\circ}$ y
ZOH casi $17^{\circ}$ (Figura~\ref{fig:disc_metodos_bode_40}). Por lo que $10$ veces es un mínimo,
y entre mayor sea la frecuencia mejores son los resultados, como se ve en la dinámica de la
Figura~\ref{fig:adrc_discreto_lazo_cerrado}.

Para el caso del ADRC basado en observador de estado extendido, debido a la relación que hay entre
las ganancias del observador $\mathbf{L}$, las del control $\mathbf{K}_p$ y los anchos de banda del
controlador y del observador, se puede ajustar su dinámica fácilmente, solo con pasar esos anchos de
banda al plano $z$ \eqref{eq:disc_ganancias_wc_wo}, por su relación con las ganancias y con la
ubicación de los polos tanto del control como del observador.

\paragraph{Rechazo de perturbaciones.} Viendo los resultados de la Figura~\ref{fig:gpi_sensibilidad},
correspondiente al Bode de la sensibilidad en lazo cerrado del GPI, vemos que podemos ajustar el
rechazo con cada ganancia (Tabla~\ref{tab:gpi_frecuencia}). A baja frecuencia,
$S_{\xi}(s) \approx (k_2/k_0)\,s$ \eqref{eq:gpi_caso_S_asintotas}, por lo que la parte lenta de la
perturbación la deciden $k_2$ y $k_0$: duplicar $k_0$ deja $k_2/k_0$ en la mitad, pero el margen de
fase baja de $65{,}8^{\circ}$ a $54{,}2^{\circ}$. El pico de $|S_{\xi}|$ lo baja sobre todo $k_1$: al
duplicarlo pasa de $40{,}3$ a $34{,}6$~dB, pero el retardo máximo baja de $110$ a $57$~s. Es decir,
lo que se gana en rechazo se paga en robustez, igual que con $\omega_o$ en el ADRC. Con la sintonía
nominal, el pico de $|S_{\xi}|$ queda en $22{,}6$~rad/h, por encima de los $12$~rad/h bajo los que está el
$95~\%$ de la energía de $\bar{\xi}$ (Tabla~\ref{tab:ident_espectro_f}), y una perturbación de $-15~\%$ baja
el oxígeno de $2{,}2$ a $2{,}15$~mg/L sin retardo y a $2{,}13$~mg/L con $\theta = 50$~s
(Figura~\ref{fig:pref_escalon}).

\paragraph{Implementaciones.} En la implementación se obtuvieron buenos resultados, logrando utilizar
un controlador de rechazo activo de perturbaciones en su forma de observador de estado y en su forma
GPI, junto con adecuaciones para su implementación, para una planta $n = 1$ y $m = 1$, por lo que
está diseñado para seguir y rechazar señales tipo paso otorgando la dinámica de una planta de orden
1. El lenguaje que se utilizó fue texto estructurado, uno de los más comunes dentro de la norma
IEC~61131-3, porque de todos es el más cercano a C. Ahora bien, a la hora de implementar, sin importar
el PLC, la marca o el lenguaje, siempre que estén dentro de la norma, se pueden llegar a implementar
estos tipos de controladores sin problemas. Solo es necesario utilizar los esquemas de las
Figuras~\ref{fig:adrc_discreto_retardos} y \ref{fig:disc_formas_directas}, que ya están en retardos,
y cualquier lenguaje de la norma, ya sea Grafcet, ladder, bloques o texto estructurado, tiene dentro
de su especificación formas de representar un retardo y hacer operaciones algebraicas. Para el ESO son
suficientes dos retardos y tres sumas, y para el GPI discreto seis retardos y un bloque de suma
(Figura~\ref{fig:disc_formas_directas}). Por lo que con cualquier enfoque se puede implementar un ADRC y con
cualquier lenguaje se pueden realizar estos retardos sin dificultad, es decir, es multiplataforma.

\paragraph{ADRC con ESO y GPI.} Al comparar los controladores con enfoque de rechazo activo de
perturbaciones, ya sea utilizando un ESO o el GPI, se ve que tienen una relación directa. En el caso
del observador de estado extendido se busca estimar la perturbación y el error de seguimiento mediante
el observador, y el GPI trata de realizar la misma idea bajo reconstrucciones de la señal, a partir de
integraciones de la señal de control y del error \eqref{eq:gpi_integrales}. Todo esto indica que son
estructuras equivalentes que tratan de realizar la misma tarea bajo enfoques diferentes. Con las
ganancias $k_0$, $k_1$ y $k_2$ queda más fácil interpretar cómo es la sensibilidad del rechazo de
perturbaciones en lazo cerrado (Figura~\ref{fig:gpi_sensibilidad}) y cómo cada ganancia puede ayudar a
modificar este efecto en el lazo, además de saber cómo afecta cada parámetro el sobrepico o el tiempo
de asentamiento (Tabla~\ref{tab:gpi_transitorio}). Mientras que el ADRC con ESO ayuda a interpretar el
rechazo activo de perturbaciones por medio de anchos de banda, donde puedo interpretar qué polos están
dedicados al rechazo y qué polos están dedicados al control y a la respuesta transitoria
\eqref{eq:sintonia_polos}.

Además, el polinomio de la sensibilidad en lazo cerrado, $e_y/\xi$, corresponde a los polos del
controlador y del observador juntos \eqref{eq:gpi_producto_S}. Por lo que vamos a tener la misma
sensibilidad, las mismas funciones de transferencia (Figura~\ref{fig:adrc_gang_kESO}), los mismos
márgenes de fase y de ganancia y los mismos anchos de banda, y lo que realmente cambia entre los
esquemas es la forma de sintonizar y su esquema de saturación. Pero incluso ahí se puede pasar de un
esquema a otro y tener el mismo resultado: con $E_1(s) = (s + \omega_o)^2$ \eqref{eq:aw_gpi_wo} el GPI
se comporta igual que el ADRC también cuando satura (Figuras~\ref{fig:sat_arranque} y
\ref{fig:gpi_discreto_vs_adrc}).

\paragraph{Comparación con el PID}
\label{sec:por_que}

Ahora bien, el ADRC también se puede escribir de una forma relacionada con el PID. En la literatura un
ADRC se puede escribir como un PID con un filtro del error, así como en este caso un PI con filtro del
error \eqref{eq:gpi_pi_filtro}. Pero su diferencia con el PID ISA con filtro en la derivada radica en
que la función de transferencia de ese PID es bipropia, y la del ADRC estrictamente propia. Esa es una
primera diferencia, donde el PID tiene peores resultados en su manejo del ruido: como vimos en la
Sección~\ref{sec:impl_ruido}, al ser bipropio deja pasar el ruido hacia la señal de control, con
$\sigma_u$ de $0{,}22~\%$ frente a $0{,}009~\%$ del ADRC (Tabla~\ref{tab:pid_escalon}), haciendo mayor
desgaste en los actuadores, a pesar de que la misma planta se comporte como un filtro.

Ahora bien, el PID se sintonizó para que su señal de control no salte más que la del ADRC y con la
misma sensibilidad máxima sin retardo, y aun así queda por detrás: es más lento, se pasa más en los
cambios de consigna y rechaza peor las perturbaciones (Tablas~\ref{tab:pid_escalon}
y \ref{tab:pid_pert_escalon}). En el PID una sola ganancia fija el rechazo, el
esfuerzo inicial y el ruido, mientras que en el ADRC el rechazo lo pone el observador. Para una planta
de primer orden como esta, un PI sintonizado con el mismo procedimiento y el mismo margen de fase del
ADRC lo supera en rechazo (Figura~\ref{fig:pid_mejores}); lo que el ADRC sigue ganando es una señal de
control que arranca suave y casi sin ruido, y la posibilidad de subir el grado del modelo de la
perturbación sin cambiar su estructura. Además, su enfoque es mucho más práctico: el seguimiento y el
rechazo se entienden y se sintonizan por separado, con dos anchos de banda, de una forma más sencilla,
simple y coherente que con el PID.

Otra diferencia está en la sintonía. Un PID también admite modificar la dinámica moviendo los polos teniendo un modelo de
la planta, pero sus ganancias $K_p$, $T_i$ y $T_d$ no se relacionan directamente con la dinámica deseada
ni con la perturbación, por lo que la mayoría de las sintonizaciones se hacen en campo, de forma
empírica. El ADRC, en cambio, se puede sintonizar por anchos de banda, para rechazar las perturbaciones
o poder ajustar una dinámica deseada del error, bajo un \emph{pole placement} con las ganancias $k_0$, $k_1$ y $k_2$,
para ajustar la función de sensibilidad en lazo cerrado y rechazar perturbaciones, o en su forma
clásica de PI más filtro.

\section{Limitaciones del estudio}
\label{sec:limitaciones}

\begin{enumerate}
  \item \textbf{No hay trazas de la planta operando con ADRC.} Dados los tiempos y las restricciones,
  es complicado tener una ventana de tiempo para realizar la implementación del código en campo dentro
  de las fechas admisibles para el proyecto. Por eso la validación es en simulación, con la planta
  identificada y la perturbación sintetizada de sus datos, y se deja todo el proyecto probado y
  verificado para su puesta en marcha cuando se tenga el espacio en campo para instalarlo.

  \item \textbf{El retardo está acotado, no medido con precisión.} El retardo se encuentra acotado y
  obtenido visualmente a partir de una de las pruebas de uno de los tres módulos dispuestos en la
  planta, por lo que dependiendo del módulo puede variar. Lo positivo de esto es que se evaluó el
  retardo con diferentes variaciones y el controlador logró responder adecuadamente, incluso mejor con
  los acondicionamientos que se hicieron.

  \item \textbf{Las variaciones se probaron una a la vez.} En las
  Figuras~\ref{fig:adrc_variacion} y \ref{fig:sat_adrc_gpi_variacion} se cambia un parámetro de la
  planta a la vez, por lo que no todas las combinaciones quedan cubiertas, y solo se validan las
  más importantes, como $b_0$ y $\theta$. Pueden aparecer casos en los que, al variar dos
  parámetros a la vez, se tengan más restricciones. Por ejemplo, si se tiene $\theta = 50$~s y al
  mismo tiempo $b_{0,\mathrm{planta}} = 0{,}327$~mg/(L$\cdot$h$\cdot$\%), un caso extremo, se tendría
  $b_{0,\mathrm{planta}}/b_0 = 2{,}18$, que supera el margen de ganancia de $2{,}12$, y el lazo
  quedaría inestable. En ese caso habría que reajustar las ganancias para aumentar el margen de
  fase y que el lazo soporte más retardo, además de suavizar la referencia con el prefiltro o con
  Bézier para disminuir el sobrepico ante el escalón. De igual forma, se recomienda que el retardo
  no sea muy alto; si lo es, toca resintonizar y, si se puede, sincronizar el observador.


  \item \textbf{Limitación del orden de la perturbación.} Se escogió para el diseño una perturbación de
  orden $m = 1$, un solo integrador, por lo que al cambiar la referencia por rampas u otros tipos se
  tendrán errores de seguimiento, relativamente pequeños en función de las ganancias de los
  controladores, pero de igual forma existirá el error de seguimiento. Para estos entornos industriales
  donde se usan PID no es tanto problema, pero si llegado el caso se necesita seguir otro tipo de
  referencias de forma exacta, $e_y(\infty) = 0$, deberá modificarse cada uno de los bloques de los
  controladores implementados y diseñados en este documento.
\end{enumerate}

\section{Implicaciones prácticas}
\label{sec:implicaciones}

Para quien tenga que decidir sobre un lazo de oxígeno en servicio, los
resultados se traducen en cinco recomendaciones concretas:

\begin{enumerate}
  \item \textbf{Medir $\bo$ en las unidades del bloque, no en unidades
  físicas.} Un solo escalón registrado a la velocidad del PLC basta, y evita
  toda la cadena de conversiones --- consigna a potencia, potencia a caudal de
  aire --- que es donde se acumula el error. Además, la ganancia se debe obtener
  siempre cerca al punto de operación, y tratar de llegar a él posteriormente con
  algún método que acondicione la referencia.

  \item \textbf{Dejar $\bo$ por encima del de la planta si hay duda.} Con $b_0$ sobreestimado la respuesta es
  más lenta, pero se conserva la robustez. Con $b_0$ subestimado el controlador es más agresivo y con
  retardo el lazo se acerca al límite de inestabilidad (Figura~\ref{fig:adrc_rlocus_b0}).

  \item \textbf{Elegir $\wo$ con el espectro de la perturbación y el ruido de
  la sonda.} Se elige como un factor de $\wc$, por encima de donde está la energía de la perturbación,
  pero sin dejar pasar el ruido de la sonda, o no lo filtraremos y tendremos una señal de control
  ruidosa. Cuando se tengan problemas de ruido, incluso es más fácil cambiar el instrumento o hacer
  mejores acondicionamientos en campo que realizar un sobreajuste en el observador.

  \item \textbf{Entregar la orden aplicada y acondicionar la referencia.} Al observador le debe entrar
  $u_{\mathrm{sat}}$ y no $u$, y los cambios de consigna deben pasar por Bézier o por el prefiltro, que
  no cambian el rechazo de perturbaciones y bajan el pico de $u$, de $93$ a $54~\%$ con
  $\theta = 50$~s (Figura~\ref{fig:pref_escalon}).

  \item \textbf{Migrar progresivamente.} Se puede arrancar el control con anchos de banda conservadores,
  no muy grandes, y dependiendo de su comportamiento en el transitorio y en el rechazo ir variando las
  ganancias, siguiendo el resumen de la Sección~\ref{sec:impl_resumen} para obtener buenos resultados.
  Como el bloque recibe las mismas señales que el PID, es una ruta de puesta en marcha que no requiere
  detener el proceso ni convencer a nadie de antemano.
\end{enumerate}
